\chapter{المتغيرات العشوائية والتوزيع الثنائي}\label{ch:g12:randvar}

يرفق \omterm{def:g12:randvar:rv}{المتغير العشوائي} عددًا بكل إمكانية من إمكانيات تجربة: ربحًا،
أو عددًا، أو مدة. وأمله الرياضي هو \omterm{def:g10:stats:mean}{متوسط} القيم التي ينتجها على
\omterm{def:g10:stats:quartiles}{المدى} الطويل، \omterm{def:g12:randvar:exp}{وتباينه} يقيس تبددها. ونجم هذا
الفصل هو \omterm{def:g12:randvar:binomial}{التوزيع الثنائي}، الذي يعدّ النجاحات في تجارب مستقلة
مكرَّرة. وقد لُقيت المتغيرات العشوائية \omterm{def:g12:randvar:binomial}{والتوزيع الثنائي}
أول مرة في \cref{ch:g11:prob,ch:g11:binom}؛ ويراجعهما هذا الفصل و
يعمّقهما، وأدوات التوافيق في \cref{ch:g12:comb} صارت الآن
متاحة.

\section{المتغيرات العشوائية المتقطعة}

\begin{definition}[المتغير العشوائي وتوزيعه]\label{def:g12:randvar:rv}
\emph{المتغير العشوائي}\index{متغير عشوائي} على مجموعة إمكانيات منتهية
$\Omega$ هو \omterm{def:g11:func:function}{دالة} $X \colon \Omega \to \R$. و
\emph{توزيعه}\index{توزيع} (أو قانونه) هو معطى قيمه الممكنة
$x_1, \dots, x_k$ والاحتمالات
\[
p_i = \P(X = x_i), \qquad \sum_{i=1}^{k} p_i = 1 .
\]
\end{definition}

\begin{definition}[الأمل الرياضي والتباين والانحراف المعياري]\label{def:g12:randvar:exp}
\emph{الأمل الرياضي}\index{أمل رياضي} للمتغير $X$ هو
\[
\E(X) = \sum_{i=1}^{k} p_i\, x_i ,
\]
و\emph{تباينه}\index{تباين} و\emph{انحرافه المعياري} هما
\[
\V(X) = \E\bigl((X - \E(X))^2\bigr) = \sum_{i=1}^k p_i\bigl(x_i - \E(X)\bigr)^2,
\qquad
\sigma(X) = \sqrt{\V(X)} .
\]
\end{definition}

\begin{proposition}[صيغة كونيغ--هويغنز]\label{prop:g12:randvar:konig}
\index{صيغة كونيغ--هويغنز}
$\V(X) = \E(X^2) - \E(X)^2$.
\end{proposition}

\begin{proof}
اكتب $m = \E(X)$ وانشر:
\[
\V(X) = \sum_i p_i (x_i^2 - 2m x_i + m^2)
= \E(X^2) - 2m\sum_i p_i x_i + m^2\sum_i p_i
= \E(X^2) - 2m^2 + m^2 . \qedhere
\]
\end{proof}

\begin{proposition}[التحويل التآلفي]\label{prop:g12:randvar:affine}
من أجل $a, b \in \R$:
\[
\E(aX + b) = a\,\E(X) + b, \qquad
\V(aX + b) = a^2\,\V(X) .
\]
\end{proposition}

\begin{proof}
الأولى إعادة \omterm{def:g10:coordgeom:system}{ترتيب} للمجموع المعرِّف. وأما الثانية، فإن $aX + b$
ينحرف عن متوسطه بمقدار $a(X - \E(X))$، والتربيع يضرب في
$a^2$.
\end{proof}

\begin{example}[الألعاب العادلة]\label{ex:g12:randvar:game}
لعبة تكلّف $2$ أورو؛ يُرمى فيها حجر نرد، ويقبض اللاعب القيمة
الظاهرة إن كانت $5$ على الأقل، ولا شيء فيما عدا ذلك. ويأخذ الربح $G$ القيم
$-2$ (\omterm{def:g10:proba:distribution}{باحتمال} $\frac46$) و $3$ ($\frac16$) و $4$ ($\frac16$):
\[
\E(G) = \frac{-8 + 3 + 4}{6} = -\frac{1}{6} < 0 .
\]
فيخسر اللاعب وسطيًا $17$ سنتًا في كل لعبة: أي إن اللعبة غير مواتية
(كما هو حال معظم الألعاب الحقيقية).
\end{example}

\section{تجارب برنولي والتوزيع الثنائي}

\begin{definition}[توزيع برنولي]\label{def:g12:randvar:bernoulli}
\emph{تجربة برنولي}\index{تجربة برنولي} هي تجربة لها إمكانيتان،
\emph{النجاح} (\omterm{def:g10:proba:distribution}{باحتمال} $p$) و\emph{الفشل}
(حيث $q = 1 - p$). ويتبع مؤشر النجاح $X$ ($X = 1$ عند النجاح و $0$ عند
الفشل) \emph{\omterm{def:g12:randvar:rv}{توزيع} برنولي} $\mathcal B(p)$:
\[
\E(X) = p, \qquad \V(X) = p(1-p) .
\]
(وفعلًا $\E(X) = p$ و $\E(X^2) = p$، وتعطي صيغة كونيغ--هويغنز أن
$\V(X) = p - p^2$.)
\end{definition}

\begin{definition}[التوزيع الثنائي]\label{def:g12:randvar:binomial}
كرّر \omterm{def:g12:randvar:bernoulli}{تجربة برنولي} $n$ مرة على نحو مستقل، وليكن $X$ العدد الكلي
للنجاحات. عندئذٍ يكون \omterm{def:g12:randvar:rv}{توزيع} $X$ هو \emph{التوزيع
الثنائي}\index{توزيع ثنائي} $\mathcal B(n, p)$.
\end{definition}

\begin{omfigure}
\begin{tikzpicture}[font=\small]
  \coordinate (root) at (0,0);
  \node (S) at (1.6,1.6) {S};
  \node (F) at (1.6,-1.6) {F};
  \node (SS) at (3.2,2.4) {S};
  \node (SF) at (3.2,0.8) {F};
  \node (FS) at (3.2,-0.8) {S};
  \node (FF) at (3.2,-2.4) {F};
  \node (SSS) at (4.8,2.8) {S};
  \node (SSF) at (4.8,2.0) {F};
  \node (SFS) at (4.8,1.2) {S};
  \node (SFF) at (4.8,0.4) {F};
  \node (FSS) at (4.8,-0.4) {S};
  \node (FSF) at (4.8,-1.2) {F};
  \node (FFS) at (4.8,-2.0) {S};
  \node (FFF) at (4.8,-2.8) {F};
  \draw[omThm, thick] (root) -- (S);
  \draw[omThm, thick] (S) -- (SS);   \draw[omThm, thick] (SS) -- (SSF);
  \draw[omThm, thick] (S) -- (SF);   \draw[omThm, thick] (SF) -- (SFS);
  \draw[omThm, thick] (root) -- (F);
  \draw[omThm, thick] (F) -- (FS);   \draw[omThm, thick] (FS) -- (FSS);
  \draw[gray] (SS) -- (SSS);
  \draw[gray] (SF) -- (SFF);
  \draw[gray] (FS) -- (FSF);
  \draw[gray] (F) -- (FF);
  \draw[gray] (FF) -- (FFS);
  \draw[gray] (FF) -- (FFF);
  \node[gray, anchor=west] at (5.3,2.8) {$X=3$};
  \node[omThm, anchor=west] at (5.3,2.0) {$X=2$};
  \node[omThm, anchor=west] at (5.3,1.2) {$X=2$};
  \node[gray, anchor=west] at (5.3,0.4) {$X=1$};
  \node[omThm, anchor=west] at (5.3,-0.4) {$X=2$};
  \node[gray, anchor=west] at (5.3,-1.2) {$X=1$};
  \node[gray, anchor=west] at (5.3,-2.0) {$X=1$};
  \node[gray, anchor=west] at (5.3,-2.8) {$X=0$};
\end{tikzpicture}

{\small ثلاث تجارب برنولي: بالضبط $\binom{3}{2} = 3$ من
المسارات $2^3$ تعطي نجاحين (بالأحمر)، كل منها \omterm{def:g10:proba:distribution}{باحتمال} $p^2(1-p)$،
إذن $\P(X = 2) = 3p^2(1-p)$.}
\end{omfigure}

\begin{theorem}\label{thm:g12:randvar:binomial}
إذا كان $X \sim \mathcal B(n, p)$، فإنه من أجل $0 \leq k \leq n$:
\[
\P(X = k) = \binom{n}{k} p^k (1-p)^{n-k},
\]
و
\[
\E(X) = np, \qquad \V(X) = np(1-p) .
\]
\end{theorem}

\begin{proof}
\omterm{def:g12:seq:sequence}{لمتتالية} محدَّدة من النتائج فيها $k$ نجاحًا و $n-k$ فشلًا
\omterm{def:g10:proba:distribution}{الاحتمال} $p^k(1-p)^{n-k}$ \omterm{def:g12:condprob:indep}{بالاستقلال}؛ وعدد هذه المتتاليات
هو عدد طرائق وضع النجاحات $k$ بين التجارب $n$،
أي $\binom nk$ (\cref{ch:g12:comb}). وبالجمع على المتتاليات نجد
الصيغة --- وتؤكد مبرهنة ثنائي الحد أن
$\sum_k \P(X = k) = (p + q)^n = 1$.

وأما \omterm{def:g12:randvar:exp}{الأمل الرياضي}، فباستعمال $k\binom nk = n\binom{n-1}{k-1}$
(\cref{exo:g12:comb:7}):
\[
\E(X) = \sum_{k=1}^{n} k\binom nk p^k q^{n-k}
= np\sum_{k=1}^{n} \binom{n-1}{k-1} p^{k-1} q^{(n-1)-(k-1)}
= np\,(p + q)^{n-1} = np .
\]
ويُبرهن على صيغة \omterm{def:g12:randvar:exp}{التباين} بالمثل بواسطة المتطابقة
$k(k-1)\binom nk = n(n-1)\binom{n-2}{k-2}$، مما يعطي
$\E(X(X-1)) = n(n-1)p^2$، ومنه
$\V(X) = n(n-1)p^2 + np - (np)^2 = np(1-p)$. (وأما البرهان البنيوي --- وهو
\omterm{def:g12:randvar:exp}{تباين} مجموع متغيرات مستقلة --- فيأتي مع
\cref{thm:g12:sums:variance}.)
\end{proof}

\begin{omfigure}
\begin{tikzpicture}
\begin{axis}[
  width=10cm, height=4.8cm,
  ybar, bar width=4pt,
  xlabel={$k$}, ylabel={$\P(X = k)$},
  xmin=-1, xmax=21, ymin=0,
  ytick={0, 0.1, 0.2},
  title={$\mathcal B(20,\ 0.3)$}]
  \addplot[fill=omDef!40, draw=omDef] coordinates {
    (0,0.000798) (1,0.006839) (2,0.027846) (3,0.071604)
    (4,0.130421) (5,0.178863) (6,0.191639) (7,0.164262)
    (8,0.114397) (9,0.065370) (10,0.030817) (11,0.012006)
    (12,0.003859) (13,0.001018) (14,0.000218) (15,0.000037)
    (16,0.000005) (17,0.000001) (18,0) (19,0) (20,0)};
\end{axis}
\end{tikzpicture}

{\small \omterm{def:g12:randvar:rv}{التوزيع} $\mathcal B(20, 0.3)$: \omterm{def:g10:stats:mean}{المتوسط} $np = 6$، \omterm{def:g11:stat:variance}{والانحراف
المعياري} $\sqrt{npq} \approx 2.05$.}
\end{omfigure}

\begin{method}[التعرف على وضعية ثنائية]\label{met:g12:randvar:recognize}
تحقق من المكونات الثلاثة قبل أن تكتب $X \sim \mathcal B(n,p)$:
\emph{عدد ثابت} $n$ من التجارب؛ و\emph{إمكانيتان} في كل تجربة \omterm{def:g10:proba:distribution}{باحتمال}
النجاح \emph{نفسه} $p$؛ و\emph{\omterm{def:g12:condprob:indep}{استقلال}} التجارب
(بالمعاينة \emph{مع} الإرجاع، أو من مجتمع كبير). ثم استعمل
\[
\P(X \geq 1) = 1 - (1-p)^n
\]
من أجل ``نجاح واحد على الأقل''، وآلة حاسبة أو جداول تراكمية من أجل
$\P(X \leq k)$ العامة.
\end{method}

\begin{example}\label{ex:g12:randvar:atleastone}
كم مرة يجب أن يُرمى حجر نرد ليكون \omterm{def:g10:proba:distribution}{احتمال} الحصول على ستة $99\%$
على الأقل؟ مع $X \sim \mathcal B\left(n, \frac16\right)$:
$\P(X \geq 1) = 1 - \left(\frac56\right)^n \geq 0.99$ يعني
$\left(\frac56\right)^n \leq 0.01$، \emph{أي}
$n \geq \frac{\ln 0.01}{\ln(5/6)} \approx 25.3$: أي ابتداءً من $n = 26$ رمية.
\end{example}

\section{تمارين}

\begin{exercise}[$\star$]\label{exo:g12:randvar:1}
\omterm{def:g12:randvar:rv}{متغير عشوائي} $X$ يأخذ القيم $-1, 0, 2, 5$ بالاحتمالات
$0.3, 0.2, 0.4, 0.1$. احسب $\E(X)$ و $\V(X)$ و $\sigma(X)$.
\end{exercise}

\begin{exercise}[$\star$]\label{exo:g12:randvar:2}
اختبار متعدد الخيارات فيه $10$ أسئلة، لكل منها $4$ خيارات، واحد
منها صحيح. ويجيب تلميذ عشوائيًا بانتظام وعلى نحو مستقل.
وليكن $X$ عدد الأجوبة الصحيحة.
\begin{enumerate}
  \item أعطِ \omterm{def:g12:randvar:rv}{توزيع} $X$ و $\E(X)$ و $\sigma(X)$.
  \item احسب $\P(X = 0)$ و $\P(X = 5)$ و $\P(X \geq 1)$.
\end{enumerate}
\end{exercise}

\begin{exercise}[$\star$]\label{exo:g12:randvar:3}
تؤمّن شركة تأمين $n = 400$ زبونًا؛ ويقدّم كل منهم مطالبة خلال
السنة \omterm{def:g10:proba:distribution}{باحتمال} $p = 0.05$، على نحو مستقل. وليكن $X$ عدد
المطالبات. حدّد \omterm{def:g12:randvar:rv}{توزيع} $X$ واحسب أمله الرياضي و
انحرافه المعياري.
\end{exercise}

\begin{exercise}[$\star\star$]\label{exo:g12:randvar:4}
في لعبة \cref{ex:g12:randvar:game}، يريد المنظّم لعبة عادلة
($\E(G) = 0$) بتغيير سعر الدخول $c$. أوجد $c$. واحسب
\omterm{def:g12:randvar:exp}{تباين} الربح في هذه النسخة العادلة؛ فهل ``العادل'' هو نفسه
``الخالي من الخطر''؟
\end{exercise}

\begin{exercise}[$\star\star$]\label{exo:g12:randvar:5}
لاعبة كرة سلة تسجّل الرميات الحرة \omterm{def:g10:proba:distribution}{باحتمال} $0.7$. وترمي
$8$ مرات (والرميات مستقلة). احسب \omterm{def:g10:proba:distribution}{احتمال} أن تسجّل:
$6$ بالضبط؛ و $6$ على الأقل؛ ومرة واحدة على الأقل. وما العدد الأرجح
من التسجيلات؟
\end{exercise}

\begin{exercise}[$\star\star$]\label{exo:g12:randvar:6}
تعرف شركة طيران أن كل مسافر حجز يحضر \omterm{def:g10:proba:distribution}{باحتمال}
$0.9$، على نحو مستقل. ولرحلة $100$ مقعد وتبيع الشركة $104$
تذكرة. عبّر، باستعمال \omterm{def:g12:randvar:binomial}{توزيع ثنائي}، عن \omterm{def:g10:proba:distribution}{احتمال} أن يحضر
مسافرون أكثر من المقاعد، وقدّر ذلك عدديًا باستعمال
آلة حاسبة (وأعطِ العبارة المضبوطة).
\end{exercise}

\begin{exercise}[$\star\star$]\label{exo:g12:randvar:7}
ليكن $X \sim \mathcal B(n, p)$. بيّن أن
\[
\frac{\P(X = k+1)}{\P(X = k)} = \frac{n-k}{k+1}\cdot\frac{p}{1-p},
\]
واستنتج أن \omterm{def:g12:randvar:rv}{التوزيع} يتزايد حتى
$k^* = \floor{(n+1)p}$ ويتناقص بعدها (وهو \emph{منوال}
\omterm{def:g12:randvar:binomial}{التوزيع الثنائي}).
\end{exercise}

\begin{exercise}[$\star\star\star$]\label{exo:g12:randvar:8}
\emph{(سان بطرسبورغ مروَّضةً.)} تُرمى قطعة نقد غير مغشوشة حتى تظهر
الصورة، لكن $10$ مرات على الأكثر. وليكن $N$ عدد الرميات المستعملة، و
يقبض اللاعب $2^N$ أورو إذا ظهرت الصورة، و $0$ فيما عدا ذلك.
\begin{enumerate}
  \item أعطِ \omterm{def:g12:randvar:rv}{توزيع} $N$ مقصورًا على الإمكانيات الرابحة:
        $\P(\text{أول صورة عند الرمية } k) = 2^{-k}$ من أجل
        $1 \leq k \leq 10$، وتحقق من \omterm{def:g10:proba:distribution}{الاحتمال} الكلي للربح.
  \item احسب المدفوع المأمول. وماذا سيصير لولا سقف
        الرميات $10$؟
\end{enumerate}
\end{exercise}

\section{مسألة: الرحلة ذات الحجز الزائد}

\begin{problem}[{مسألة نهاية الأسبوع --- شركات الطيران تبيع مقاعد أكثر مما
تملك، والمصانع تقبل دفعات بالكاد فحصتها، و
التوزيع الثنائي يحكّم بينهما}]
\label{pb:g12:randvar:1}
شركة طيران فيها $100$ مقعد تبيع بسرور $105$ تذكرة: فنحو
$10\,\%$ من المسافرين لا يحضرون أبدًا، والمقاعد الفارغة لا تكسب
شيئًا. فإلى أي حد تستطيع الشركة أن تدفع قبل أن يأكل المسافرون المُنزَلون
الربح؟ يجيب \omterm{def:g12:randvar:binomial}{التوزيع الثنائي}
(\cref{thm:g12:randvar:binomial}) بدقة الرقم العشري --- و
الآلة نفسها تفحص دفعات المصانع، وتسعّر اليانصيبات، و
تبقي المؤمّنين قادرين على الدفع. قرارات تحت التكرار: هذه هي
وظيفة \omterm{def:g12:randvar:binomial}{التوزيع الثنائي} اليومية.

\medskip
\noindent\textbf{الجزء الأول --- التمكّن.}
\begin{enumerate}
  \item أعطِ جدول \omterm{def:g12:randvar:rv}{التوزيع} الكامل لعدد $X$ من
        الأوجه الستة في ثلاث رميات لحجر نرد.
  \item احسب $\E(X)$ و $V(X)$
        (\cref{def:g12:randvar:exp}).
  \item كشك يتقاضى $1$ أورو مقابل ثلاث رميات ويدفع $2$
        أورو عن كل ستة تظهر. احسب الربح المأمول: فهل اللعبة
        عادلة؟
  \item تدرَّب على قائمة التحقق (\cref{met:g12:randvar:recognize}):
        هل عدد القلوب في $5$ أوراق تُسحب
        \emph{بدون} إرجاع ثنائي؟ وماذا لو كان مع الإرجاع؟
        برّر.
  \item من أجل $X \sim \mathcal B(20,\ 0.3)$: أعطِ $\E(X)$ و
        $\sigma(X)$، ثم $\E(S)$ و $\sigma(S)$ من أجل النتيجة
        $S = 5X - 10$ (\cref{prop:g12:randvar:affine}).
\end{enumerate}

\medskip
\noindent\textbf{الجزء الثاني --- الرحلة ذات الحجز الزائد.}
المقاعد: $100$. والتذاكر المبيعة: $105$. ويحضر كل حامل تذكرة
\omterm{def:g10:proba:distribution}{باحتمال} $0.9$، على نحو مستقل؛ وليكن $X$ عدد الحاضرين.
\begin{enumerate}[resume]
  \item النمذجة: برّر $X \sim \mathcal B(105,\ 0.9)$ بقائمة
        التحقق --- واعترف بأضعف نقطة في النموذج
        (فهل عدم الحضور مستقل فعلًا؟ فكّر في المجموعات
        والعواصف).
  \item احسب $\E(X)$ و $\sigma(X)$.
  \item يحدث الإنزال عندما يكون $X \geq 101$: اكتب
        $\P(X \geq 101)$ مجموعًا صريحًا من خمسة حدود
        ثنائية.
  \item قدّر المجموع (بآلة حاسبة): فما نسبة الرحلات
        التي ترى مسافرًا مُنزَلًا واحدًا على الأقل؟
  \item المال: تجلب التذاكر الخمس الزائدة
        $5 \times 200 = 1\,000$ أورو؛ ويكلّف كل مسافر مُنزَل
        $800$ أورو تعويضًا. احسب العدد المأمول
        للمسافرين المُنزَلين،
        $\sum_{k=101}^{105} (k - 100)\,\P(X = k)$، والتعويض المأمول،
        ثم الحكم على السياسة.
  \item التصميم: يتسامح المنظّم مع
        $\P(\text{الإنزال}) < 5\,\%$. وباختبار
        $n = 105, 106, 107, \dots$ تذكرة، أوجد أكبر $n$
        مسموح به.
  \item اختصار الجرس: من أجل $n = 105$، احسب \omterm{pb:g11:stat:1}{الدرجة المعيارية}
        لعتبة الإنزال،
        $z = \frac{100.5 - \E(X)}{\sigma(X)}$، وقارن التقدير
        الخام ``نحو $2\sigma$، أي نحو $2$--$3\,\%$ في
        الذيل الأعلى'' بجوابك المضبوط. (\omterm{def:g10:functions:graph}{والمنحنى}
        وراء هذا الاختصار هو نجم الفصول
        التالية.)
\end{enumerate}

\medskip
\noindent\textbf{الجزء الثالث --- بوابة المصنع.}
تُقبل دفعة من القطع إذا احتوت \omterm{def:g10:proba:sample}{عينة} من $20$ قطعة على قطعة
معيبة واحدة على الأكثر.
\begin{enumerate}[resume]
  \item إذا كان معدل العيب الحقيقي $2\,\%$ (وهي دفعة نزيهة)،
        فاحسب \omterm{def:g10:proba:distribution}{احتمال} القبول.
  \item وإذا كان المعدل $10\,\%$ (وهي دفعة رديئة)، فاحسبه من جديد.
  \item سمِّ خطري الإجراء (الدفعة النزيهة
        المرفوضة؛ والدفعة الرديئة المقبولة)، واقرأ قيمتيهما من
        السؤالين 13--14، وقل ما التغيير الوحيد الذي يحسّن
        كليهما معًا --- وبأي ثمن.
  \item وراء التحسين: بيّن أن \emph{\omterm{def:g10:stats:series}{تواتر}} العيب
        الملاحَظ $\frac Xn$ في \omterm{def:g10:proba:sample}{عينة} حجمها $n$ له
        \omterm{def:g11:stat:variance}{انحراف معياري} $\sqrt{\frac{p(1-p)}{n}}$، ثم
        اختم بكيفية تحسّن الدقة مع حجم \omterm{def:g10:proba:sample}{العينة} (وهو صديق
        قديم: $\frac{1}{\sqrt n}$ الذي رأيناه من قبل).
\end{enumerate}

\medskip
\noindent\textbf{الجزء الرابع --- الألعاب واليانصيبات والاحتياطيات.}
\begin{enumerate}[resume]
  \item لعبة سان بطرسبورغ المسقوفة في
        \cref{exo:g12:randvar:8} مدفوعها المأمول $10$
        أورو. قارنها في فقرة واحدة بالمفارقة غير المسقوفة
        التي لقيناها في \cref{pb:g11:prob:1}: فما الذي يروّضه
        السقف بالضبط، وما السعر العادل الآن؟
  \item يانصيب خيري يبيع $200$ تذكرة بسعر $2$ أورو؛ و
        الجوائز: سلة قيمتها $100$ أورو وسلتان قيمة كل منهما $50$ أورو.
        احسب الربح المأمول للتذكرة وهامش اليانصيب؛
        وقارنه بهامش بيت الروليت الأوروبي
        البالغ $\frac{1}{37} \approx 2.7\,\%$. فلماذا يفلت اليانصيب
        بفعلته؟
  \item مؤمّن يحمل $1\,000$ عقد مستقل: \omterm{def:g10:proba:distribution}{احتمال} المطالبة
        $0.01$ لكل منها، وحجم المطالبة $10\,000$ أورو، و
        القسط $130$ أورو. احسب الربح السنوي المأمول
        \emph{و} \omterm{def:g11:stat:variance}{الانحراف المعياري} لمجموع
        المطالبات. وقارن العددين: فما الذي تفرضه المقارنة
        على المؤمّنين الحقيقيين أن يحتفظوا به؟
  \item الخاتمة --- \omterm{def:g12:randvar:binomial}{التوزيع الثنائي} حكمًا للقرارات:
        قائمة التعرف؛ والبوصلة $\E \pm \sigma$؛
        واحتمالات الذيول ثمنًا للخطر؛ والتحذيران القائمان
        (\omterm{def:g12:condprob:indep}{الاستقلال} ادعاء نمذجة، وأن الذيول
        النادرة، لا المتوسطات، هي التي تسبّب الخراب). جملة واحدة لكل بند،
        مع الإشارة: قانون الأعداد الكبيرة \omterm{def:g10:functions:graph}{ومنحنى}
        الجرس، في الفصول التالية، يكملان كتاب قواعد الحكم.
\end{enumerate}
\end{problem}
